% =============================================================================
% Chapter 15: Results -- Part III (Extended Validation)
% =============================================================================
\mychapter{15. Results --- Extended Validation}

\mysection{15.1 Loss-Function Ablation Results}
Table~\ref{tab:loss_ablation} reports validation-set performance for all four loss configurations,
trained and evaluated under the identical split described in Section~14.3.

\begin{table}[htbp]
\centering
\caption{Loss-function ablation, validation-sample (sample\_03) results. BCE achieves the best raw
  Dice and IoU of the four configurations.}
\label{tab:loss_ablation}
\begin{tabular}{lllll}
\toprule
\textbf{Loss} & \textbf{Val Dice} & \textbf{Val IoU} & \textbf{Precision} & \textbf{Recall} \\
\midrule
BCE                    & \textbf{0.7275} & \textbf{0.6235} & 0.809 & 0.7206 \\
Dice                   & 0.691           & 0.578           & 0.822 & 0.672 \\
Combined Dice-Focal (Part II's choice) & 0.696 & 0.578      & 0.741 & 0.7278 \\
Focal                  & 0.574           & 0.466           & \textbf{0.976} & 0.471 \\
\bottomrule
\end{tabular}
\end{table}

BCE achieves the best Dice and IoU of the four configurations. It was not selected as the primary
formulation when this comparison was first run, combined Dice-Focal was retained instead on the
grounds of a more balanced precision/recall trade-off, but re-examining the exact figures shows that
framing overstated the difference: BCE's recall (0.7206) is only 0.0072 below combined Dice-Focal's
(0.7278), not a meaningful trade-off given the other three metrics all favour BCE outright. Separately,
and independently of this ablation, the held-out generalisation-gap checkpoint identified in
Section~15.2 below records a validation Dice of 0.7275 at epoch 13, an exact match to this table's BCE
row, confirming that checkpoint was already trained with BCE loss. Taken together, these two points of
evidence support formalising BCE as the primary loss function for this pipeline going forward, rather
than combined Dice-Focal.

\mysection{15.2 The Generalisation Gap}
Table~\ref{tab:generalisation_gap} reports the same checkpoints' Dice on their validation sample
(sample\_03) against their Dice on the fully held-out test sample (sample\_01, never included in
training or validation for any configuration in this part).

\begin{table}[htbp]
\centering
\caption{Validation Dice vs.\ held-out test Dice, across the preprocessing and masking ablations of
  Sections~14.3--14.5. All configurations share the same held-out sample, sample\_01. Acceptance
  criteria (Dice~$\geq$~0.75, IoU~$\geq$~0.60, Recall~$\geq$~0.80) fail in every configuration on the
  held-out sample.}
\label{tab:generalisation_gap}
\small
\begin{tabular}{p{7.8cm}p{3.2cm}p{3.2cm}}
\toprule
\textbf{Configuration} & \textbf{Validation Dice} & \textbf{Held-out test Dice} \\
\midrule
BCE, baseline (3-stage) preprocessing, circular mask & 0.7275 & 0.1218 \\
BCE, + BHC/ring preprocessing, circular mask         & 0.7013 & 0.1693 \\
BCE, + BHC/ring preprocessing, shape-aware mask       & $\sim$0.70 & 0.2119 \\
\bottomrule
\end{tabular}
\end{table}

Every configuration shows the same pattern: validation Dice in the 0.70--0.73 range, held-out test
Dice far lower, 0.12--0.21. This gap is present in the original (Part II) configuration as well
(Section~10.3: validation Dice 0.600, held-out test Dice 0.147), so it is corroborated by two
independently run experiments using different sample-role assignments and different loss/preprocessing
configurations, not an artefact of one specific run. Table~\ref{tab:preproc_ablation} reports the full
held-out test metrics for the two directly comparable configurations (baseline vs.\ BHC/ring,
circular mask both).

\begin{table}[htbp]
\centering
\caption{Preprocessing ablation, full held-out test set metrics (sample\_01), circular mask both
  configurations.}
\label{tab:preproc_ablation}
\begin{tabular}{lllll}
\toprule
\textbf{Configuration} & \textbf{Dice} & \textbf{IoU} & \textbf{Precision} & \textbf{Recall} \\
\midrule
Baseline (3-stage preprocessing) & 0.1218 & 0.0907 & 0.2908 & 0.3134 \\
+ BHC and ring-artefact correction & 0.1693 & 0.1328 & 0.4049 & 0.2974 \\
\bottomrule
\end{tabular}
\end{table}

Held-out test Dice improves by 39\% relative (0.1218 $\to$ 0.1693) with BHC/ring correction added,
while validation Dice moves slightly in the opposite direction (0.7275 $\to$ 0.7013). Read together,
this indicates the preprocessing change improves genuine generalisation to unseen data rather than
fitting the validation sample more closely, since the more trustworthy of the two numbers (held-out,
never used for any training decision) improved while the less trustworthy one (validation, used for
checkpoint selection) did not. Adding the shape-aware boundary mask of Section~14.5 on top of BHC/ring
continues the same trend, held-out Dice rises further to 0.2119, a monotonic improvement across all
three configurations (0.1218 $\to$ 0.1693 $\to$ 0.2119) even though every configuration still fails the
thesis's acceptance criteria by a wide margin.

\mysection{15.3 Shape-Aware Masking Results}
Table~\ref{tab:podfam_shape_aware} reports the corrected full-volume Bernsen defect statistics for
sample\_06 and sample\_07 using the shape-aware boundary method of Section~14.5, alongside the
uncorrected result the original erosion setting produced.

\begin{table}[htbp]
\centering
\caption{sample\_06/07, shape-aware boundary detection, before and after erosion-margin correction.
  sample\_06's initial result was found to be a boundary artefact, not genuine porosity; sample\_07's
  corrected result was confirmed genuine by visual inspection of the full slice mosaic.}
\label{tab:podfam_shape_aware}
\small
\begin{tabular}{p{2.2cm}p{6.5cm}p{5.5cm}}
\toprule
\textbf{Sample} & \textbf{Defect fraction, erosion=8px (uncorrected)} & \textbf{Defect fraction, erosion=35px (corrected)} \\
\midrule
sample\_06 & 2.92\% (boundary rim artefact) & \textbf{0.030\%} \\
sample\_07 & --- & \textbf{0.759\%} \\
\bottomrule
\end{tabular}
\end{table}

sample\_06's uncorrected result (2.92\%) is roughly two orders of magnitude above its corrected value
(0.030\%), traced to a thin rim of partial-volume boundary pixels being misclassified as defect at the
original 8~px erosion margin; a sweep of erosion values (25, 30, 35, and 40~px) showed the false rim
shrinking smoothly and predictably as the margin increased, with 35~px selected as the value fully
removing the artefact without eroding genuine interior material. sample\_07's corrected result
(0.759\%) was cross-checked by direct visual inspection of its full slice mosaic and judged to reflect
genuine porosity rather than a residual boundary effect. For comparison, the circular-fit method
applied to sample\_01--05 in this same later pipeline pass gives full-volume Bernsen defect fractions
of 0.055\% (sample\_01), 6.494\% (sample\_02), 0.181\% (sample\_03), 1.641\% (sample\_04, adaptive
Otsu boundary threshold), and 53.42\% (sample\_05, fixed boundary threshold, required because its
$\sim$72\% porosity leaves too few solid pixels for Otsu's two-class assumption to hold, Section~14.5).
These figures are from an independently re-run pipeline pass and are not identical to Section~10.1's
Table~\ref{tab:podfam_classical} figures for the same samples, both are genuine pipeline outputs from
different points in this work and are reported here as such rather than reconciled into one number.

The Pearson correlation between each classical method's per-sample defect fraction and Kim et
al.'s~\citep{Kim2017} independently published porosity (Section~14.9, five NIST samples in common)
was recomputed under the circular, no-mask, and shape-aware boundary configurations, to check whether
replacing the circular assumption changes the pipeline's already-established external-validation
accuracy. Table~\ref{tab:mask_r_comparison} shows the result.

\begin{table}[htbp]
\centering
\caption{Pearson $r$ vs.\ Kim et al.'s published porosity, by boundary-detection method (NIST samples,
  $n=5$). Removing the boundary mask entirely is unusable; shape-aware masking is statistically
  indistinguishable from the circular mask it replaces.}
\label{tab:mask_r_comparison}
\begin{tabular}{llll}
\toprule
\textbf{Method} & \textbf{Circular mask} & \textbf{No mask} & \textbf{Shape-aware mask} \\
\midrule
Otsu    & 0.983 & 0.984 & 0.982 \\
Yen     & 0.503 & $\sim$0.43 & $\sim$0.50 \\
Bernsen & \textbf{0.995} & 0.992 & 0.991 \\
\bottomrule
\end{tabular}
\end{table}

With no boundary mask at all, background air is misclassified as defect on low-porosity samples, Otsu's
defect fraction on one near-zero-porosity sample jumped from 0\% to 24\% once the mask was removed,
making that configuration unusable despite its $r$ value alone looking comparable. Shape-aware masking,
by contrast, reproduces the circular mask's accuracy to within 0.001--0.004 of $r$ across all three
methods, recovering the same measured accuracy with a method that carries no assumption about
cross-sectional shape.

\mysection{15.4 Background-Masking Fix for Detector Inference}
Table~\ref{tab:bg_mask_fix} reports the trained detector's predicted defect fraction on sample\_07
(out-of-distribution inference only, Section~14.7) before and after applying the shape-aware
background-masking fix.

\begin{table}[htbp]
\centering
\caption{sample\_07, U-Net predicted defect fraction, same checkpoint and same sample, before and
  after the background-masking fix.}
\label{tab:bg_mask_fix}
\begin{tabular}{ll}
\toprule
\textbf{Configuration} & \textbf{Predicted defect fraction} \\
\midrule
2.5D U-Net, before fix    & 26.778\% \\
2.5D U-Net, after fix     & 4.764\% \\
2D U-Net, after fix       & 2.691\% \\
Bernsen classical (reference, same sample) & 7.060\% \\
\bottomrule
\end{tabular}
\end{table}

The fix produces a 5.6$\times$ reduction (26.778\% $\to$ 4.764\%) using the identical checkpoint on the
identical sample, only the background masking changed. This confirms most of the pre-fix number was
background hallucination on transition/support-structure content rather than a fundamental defect-
detection failure. sample\_07 was never in this checkpoint's training distribution (different
geometry, and processed with a different preprocessing configuration than the training samples), so
none of the figures in this table can be checked against ground truth; they are reported as an
internal consistency check on the fix, not as a validated accuracy result.

\mysection{15.5 Canny Boundary Detection: Validation and Failure}
On sample\_01--05, Canny edge detection with morphological closing reproduces the production
circular-fit defect fractions closely, 6.62\% vs.\ 6.49\% for sample\_02, and 1.69\% vs.\ 1.64\% for
sample\_04, cross-confirming those production numbers by an independent method. On sample\_06/07, it
correctly follows the true material boundary through most of each volume's depth, but on two slices it
instead traces the surrounding support scaffold, reporting 29.96\% defect content against a true value
of approximately 0--2\% at those slices. As discussed in Section~14.6, this failure was traced to the
underlying image content (a genuinely ambiguous material/scaffold contrast at those specific slices)
rather than a fixable parameter, and the shape-aware method of Section~14.5 remains the production
choice for these two samples for that reason.

\mysection{15.6 Conv-Block Attribution Results}
Table~\ref{tab:convblock} reports each convolutional block's relative contribution to the defect
decision, computed as described in Section~14.8, on the BCE checkpoint identified in Section~15.1 as
the strongest of the four loss-function configurations.

\begin{table}[htbp]
\centering
\caption{Gradient$\times$activation attribution per convolutional block, BCE checkpoint (verified
  val\_dice = 0.7275, epoch 13). The three encoder blocks together account for approximately 69\% of
  the total attributed signal; the final decoder block is the least influential of all nine.}
\label{tab:convblock}
\begin{tabular}{lll}
\toprule
\textbf{Block} & \textbf{Channels} & \textbf{Attributed signal} \\
\midrule
encoder\_1  & 64   & 27.78\% \\
encoder\_2  & 128  & 25.81\% \\
encoder\_3  & 256  & 15.29\% \\
encoder\_4  & 512  & 6.84\% \\
decoder\_3  & 128  & 6.36\% \\
decoder\_2  & 256  & 6.10\% \\
bottleneck  & 1024 & 5.92\% \\
decoder\_1  & 512  & 4.63\% \\
decoder\_4  & 64   & \textbf{1.28\%} \\
\bottomrule
\end{tabular}
\end{table}

encoder\_1 through encoder\_3 together account for 68.88\% of the attributed signal, while decoder\_4,
the block closest to the final output layer, is the least influential of all nine (1.28\%). This
analysis was first run, in error, on the combined Dice-Focal checkpoint (\texttt{top6\_by\_voids})
rather than the correct BCE checkpoint, before the ablation of Section~15.1 had settled which
checkpoint should be treated as primary; that run found the same qualitative pattern at a smaller
magnitude (encoder\_1--3 combined: 61.7\%; decoder\_4: 1.46\%). Re-running the analysis on the verified
BCE checkpoint after the fact, the pattern reported in Table~\ref{tab:convblock} replicated cleanly and
was, if anything, stronger.

\mysection{15.7 External Validation and Voxel-Size Results}
The gravimetric cross-check identified in Section~14.9 covers, in Kim et al.'s own sample numbering,
samples 2, 4, 5, and 6, with XCT-vs-gravimetric agreement within 1--4\% absolute and threshold-
sensitivity uncertainty below $\pm 0.12\%$~\citep{Kim2017}. This strengthens the interpretation of this
thesis's own $r = 0.995$ Bernsen-vs-Kim-et-al.\ correlation, Table~\ref{tab:mask_r_comparison}, as
agreement with a reference value that already carries independent physical grounding rather than a
purely XCT-internal comparison, with the caveat, stated in Section~14.9, that individual per-sample
confirmation of the gravimetric check for all five NIST samples used in this thesis's own correlation
was not independently verified here.

The physical voxel-size figures determined in Section~14.10 are: PODFAM sample\_06 and sample\_07,
12.001~\textmu m/px (from embedded raw-TIFF resolution metadata, consistent across both samples),
giving approximate physical dimensions of 12.3~mm~$\times$~14.8~mm~$\times$~17.6~mm for sample\_06 and
14.6~mm~$\times$~13.7~mm~$\times$~17.7~mm for sample\_07; and NIST CoCr samples, approximately
2.5~\textmu m/px for samples 1, 2, and 4, and 0.87~\textmu m/px for samples 3 and 5, per the dataset's
own published documentation~\citep{NIST_CoCr}. \texttt{config.py}'s \param{PIXEL\_SIZE\_UM} constant
itself remains at its unfilled placeholder value in the running pipeline; these figures are reported
as an external measurement against the raw data and documentation, not as a change already applied to
the pipeline's own unit conversions.
